\section{Conclusion}
\label{sec:conclusion}

The case reference house sets a physical limit: with unlimited work, four abstract servers could
deliver \qty{18129}{kWh} of useful heat, 58\% of the demand, and a simple local controller already
reaches the full-year upper bound. With finite, sourced job demand the picture changes. The same
four servers deliver \qty{7133}{kWh} (22.8\%), dump nothing, and miss 8.9\% of deadlines, nearly
all for work that arrives in summer when the house only needs hot water. Demand, not heat, limits
the house in winter; heat limits it in summer.

Economic viability is an actor-level question. Income exceeds the electricity cost, but at the sourced demand four GPU servers
lose about CHF~14,300 a year after hardware. One server earns almost the same income and is the
only configuration that covers its hardware: the combined system keeps about CHF~5,250 a year, the
operator alone, who gives the heat away, about CHF~4,470 (payback 3.4 years). Four servers need 2.8 times the demand or 2.6 times
the Type~1 price; raising the research price alone does not suffice. The sizing question of the
case therefore has a practical answer under these assumptions: size for the work that can be sold,
not for the winter peak.

Assuming 86\% boiler efficiency, the case formula describes a real house's space heating within
1.3\% of the case's design power,
while that house's summer heat use is three times the case's hot water and removes most summer
deadline misses. Year-round heat sinks serve the same role in a network: a nearby pool takes 76\%
of the July work, and in synthetic experiments with earlier research strategies, routing reduces missed deadlines where houses alone
cannot. The demonstrator shows the scheduler freezing and resuming a real container on the heat
signal, with measured CPU time and energy.

The next tests are a measured, thermally integrated node, a customer commitment stating price and
minimum paid capacity, and a measured hot-water profile; the thermal test must record delivery
temperature, circulation-pump power and storage losses so that measured useful heat can replace
the idealised recovery assumption. The paper's contribution is a reproducible boundary between
demonstrated heat control, conditional home economics and network scenarios.
